\documentclass[11pt]{article}
\usepackage[margin=0.95in]{geometry}
\usepackage{amsmath,amssymb,mathtools,bm}
\usepackage{xcolor}
\usepackage[hidelinks]{hyperref}
\usepackage[capitalise]{cleveref}

\newcommand{\wt}{\widetilde}
\newcommand{\Original}[1]{\textcolor{red!70!black}{#1}}
\newcommand{\Corrected}[1]{\textcolor{blue!65!black}{#1}}
\newcommand{\Issue}[1]{\par\medskip\noindent\textbf{Issue.} #1\par}
\newcommand{\Fix}[1]{\par\noindent\textbf{Correction.} #1\par}
\newcommand{\Reason}[1]{\par\noindent\textbf{Reason.} #1\par\medskip}
\newcommand{\dd}{\mathrm{d}}

\title{Marked Changes, Audit Trail, and Referee-Oriented Notes\\
for the ring-with-shortcut first-passage calculation}
\author{}
\date{\today}

\begin{document}
\maketitle
\tableofcontents

\section{How to read this audit file}

This file is not merely a shorter duplicate of the corrected derivation.  It records what changed,
why it changed, which parts of the original calculation can be retained, which parts should be
replaced, and what the corrected analytical answer says about the bimodality/decay question.  Red
formulae labelled ``original'' are formulae or logical steps that should not be used in the revised
argument under the stochastic interpretation stated below.  Blue formulae labelled ``corrected'' are
replacement formulae.  The accompanying file
\texttt{Calculations\_corrected\_complete\_v3.tex} contains the fully expanded derivation with the
intermediate algebra left in.

The stochastic interpretation used for the correction is the one that appears to be intended by the
wording ``a shortcut from $u$ to the absorbing site $v$'': at a visit to $u$, a probability
\begin{equation}
  \lambda=\beta(1-q)
\end{equation}
is diverted to the target $v$.  Since $v$ is absorbing, this diverted probability exits the
transient state space immediately.  If the intended model is instead a probability-conserving
shortcut on a ring before imposing absorption, or a reverse shortcut, then the model statement must
be rewritten explicitly because the algebra describes a different Markov chain.

\section{Executive audit result}

The absorbing-ring calculation is mostly salvageable, but it needs several local corrections:
index ranges, odd/even mode notation, a sign convention in the double sum, corrected orthogonality
checks, and several typos.  The shortcut-to-target calculation requires a substantive correction.
The original shortcut section uses the denominator
\begin{equation}
  \Original{1-z\beta(1-q)\wt W_u(u,z)}
\end{equation}
and consequently introduces
\begin{equation}
  \Original{
  \wt H_{N,N/2}(z)=\frac{1}{\dfrac{q}{\beta(1-q)}-\phi(z)}.}
\end{equation}
Under the stochastic shortcut-to-absorbing-target interpretation, the transient dynamics instead
contains an additional killing channel at $u$, and the corrected denominator is
\begin{equation}
  \Corrected{1+z\beta(1-q)\wt W_u(u,z)}.
\end{equation}
Thus the symmetric-case auxiliary factor is
\begin{equation}
  \Corrected{
  \wt H^{\rm corr}_{N,N/2}(z)=\frac{1}{\dfrac{q}{\beta(1-q)}+\phi(z)}.}
\end{equation}
Equivalently, if the same plotted function $\phi(1/s)$ is retained, the horizontal line must be
$-q/[\beta(1-q)]$, not $+q/[\beta(1-q)]$.

Therefore the original threshold
\begin{equation}
  \Original{\beta\leq \frac{2q}{(1-q)N}}
\end{equation}
is not a valid analytical condition for the corrected stochastic model.  The corrected long-time
right-tail decay is
\begin{equation}
  \Corrected{F_\rho(t)\sim B_{\rho 1}s_1^{t-1}},
\end{equation}
where $s_1$ is determined implicitly by
\begin{equation}
  \Corrected{aT_L(y_s)+U_{L-1}(y_s)=0,
  \qquad a=\frac{q}{\beta(1-q)},\qquad y_s=1+\frac{s-1}{q}.}
\end{equation}
It is not generically $t\alpha_1^t$.  For weak shortcut strength,
\begin{equation}
  \Corrected{s_1=\alpha_1-\frac{2\beta(1-q)}{N}+O([\beta(1-q)]^2),}
\end{equation}
so the corrected decay gap is
\begin{equation}
  \Corrected{1-s_1=1-\alpha_1+\frac{2\beta(1-q)}{N}+O([\beta(1-q)]^2).}
\end{equation}
For large $N$,
\begin{equation}
  \Corrected{1-s_1\sim \frac{q\pi^2}{2N^2}+\frac{2\beta(1-q)}{N}.}
\end{equation}
This last line is the cleanest referee-facing decay-dependence statement.

\section{Audit item 1: the model must be stated before using the defect formula}

\Issue{The original derivation starts from a shortcut defect formula for a probability-conserving
ring propagator and then states that one may replace the probability-conserving propagator by the
absorbing propagator.  This replacement is not innocuous when the shortcut endpoint is itself the
absorbing site.}

\Fix{State the transient-space transition matrix explicitly.  Let $P_0$ be the transition matrix
for the ring with absorbing site $v$ removed from the transient state space.  If the shortcut from
$u$ to $v$ replaces probability $\lambda=\beta(1-q)$ of staying at $u$, then the transient matrix is
\begin{equation}
  P=P_0-\lambda |u\rangle\langle u|.
\end{equation}
The sign is negative because probability leaves the transient set at $u$.}

\Reason{The resolvent is
\begin{align}
  (I-zP)^{-1}
  &=\bigl(I-zP_0+z\lambda |u\rangle\langle u|\bigr)^{-1}.
\end{align}
Using the rank-one Sherman--Morrison identity gives
\begin{equation}
  \Corrected{
  \wt S_{n_0}(n,z)=\wt W_{n_0}(n,z)
  -\frac{z\lambda\wt W_{n_0}(u,z)\wt W_u(n,z)}{1+z\lambda\wt W_u(u,z)}.}
\end{equation}
The original plus correction would increase transient occupation, which is the opposite of a
shortcut into an absorbing target.}

\section{Audit item 2: corrected first-passage generating function}

\Issue{The original first-passage generating function has the shortcut correction with the wrong
sign inherited from the transient propagator.}

\Fix{For the killed process, survival and first-passage are related by
\begin{equation}
  \wt R_{n_0}(z)=\sum_{n\neq v}\wt S_{n_0}(n,z)=\frac{1-\wt F_{n_0}(v,z)}{1-z}.
\end{equation}
Using the corrected $\wt S$ and the identity
\begin{equation}
  \sum_{n\neq v}\wt W_a(n,z)=\frac{1-\wt F_a^{(0)}(v,z)}{1-z},
\end{equation}
we obtain
\begin{equation}
  \Corrected{
  \wt F_{n_0}(v,z)=\wt F^{(0)}_{n_0}(v,z)+
  \frac{z\lambda\wt W_{n_0}(u,z)[1-\wt F_u^{(0)}(v,z)]}
  {1+z\lambda\wt W_u(u,z)}.}
\end{equation}}

\Reason{This correction is positive for $0\leq z<1$ because $\wt W$ and $1-\wt F$ have
non-negative coefficients.  It passes the probabilistic monotonicity check: adding an absorbing
shortcut cannot postpone absorption in the stochastic coupling where paths are identical until the
first time a shortcut event occurs.}

\section{Audit item 3: symmetric case and replacement for the original Fig. 1 analysis}

Let $N=2L$, $v=0$, $u=L$, and let $\rho$ denote the starting distance from the target along either
half-ring.  Define
\begin{equation}
  a=\frac{q}{\lambda}=\frac{q}{\beta(1-q)},
  \qquad
  y=1+\frac{1}{q}\left(\frac1z-1\right).
\end{equation}
The fully expanded derivation gives
\begin{equation}
  \Corrected{
  \wt F_\rho(z)=
  \frac{aT_{L-\rho}(y)+U_{\rho-1}(y)+U_{L-\rho-1}(y)}
  {aT_L(y)+U_{L-1}(y)}.}
  \label{eq:marked-main-F}
\end{equation}
This formula should replace the original pole analysis based on
\begin{equation}
  \Original{\frac{1}{a-\phi(z)}}.
\end{equation}
The denominator of \eqref{eq:marked-main-F} is
\begin{equation}
  D(y)=aT_L(y)+U_{L-1}(y).
\end{equation}
Using the Chebyshev identity
\begin{equation}
  T_L(y)=\frac12\left[U_L(y)-U_{L-2}(y)\right]
\end{equation}
and, for $N=2L$,
\begin{equation}
  U_{2L-1}(y)=2T_L(y)U_{L-1}(y),
\end{equation}
we can also write the corrected denominator as a shifted version of the original plotted object:
\begin{equation}
  a+\frac{U_{L-1}(y)}{T_L(y)}=a+\phi(z).
\end{equation}
Thus the same curve $\phi(1/s)$ may be kept only if the horizontal line is moved to $-a$.

\section{Audit item 4: corrected pole ordering and the missing physical check}

Set $y_s=1+(s-1)/q$.  The corrected pole equation is
\begin{equation}
  aT_L(y_s)+U_{L-1}(y_s)=0.
\end{equation}
For the largest root put $y_s=\cos\theta$.  Since
\begin{equation}
  T_L(\cos\theta)=\cos(L\theta),
  \qquad
  U_{L-1}(\cos\theta)=\frac{\sin(L\theta)}{\sin\theta},
\end{equation}
the pole equation is
\begin{equation}
  a\cos(L\theta)+\frac{\sin(L\theta)}{\sin\theta}=0,
  \qquad\text{or}\qquad
  \Corrected{\tan(L\theta)=-a\sin\theta.}
\end{equation}
The largest root lies in
\begin{equation}
  \frac{\pi}{L}<L\theta_1<\frac{\pi}{2}
  \quad\text{is false, whereas}\quad
  \Corrected{\frac{\pi}{2}<L\theta_1<\pi}
\end{equation}
for the first interval relevant to the largest positive $s$.  Equivalently,
\begin{equation}
  \Corrected{\frac{\pi}{2L}<\theta_1<\frac{\pi}{L}.}
\end{equation}
Therefore
\begin{equation}
  \Corrected{1-q+q\cos\frac{\pi}{L}<s_1<1-q+q\cos\frac{\pi}{2L},}
\end{equation}
that is
\begin{equation}
  \Corrected{\gamma_1<s_1<\alpha_1,}
  \qquad
  \gamma_1=1-q+q\cos\frac{2\pi}{N},
  \qquad
  \alpha_1=1-q+q\cos\frac{\pi}{N}.
\end{equation}
This inequality is a necessary physical sanity check.  A shortcut into the target must accelerate
the ultimate tail relative to the no-shortcut absorbing ring; hence $s_1$ should be below
$\alpha_1$, not above it.

\section{Audit item 5: the beta condition and the right-tail decay}

\Issue{The original conclusion compares $s_1$ with $\alpha_1$ using the incorrect Fig. 1 geometry
and obtains a threshold $\beta\le 2q/[(1-q)N]$.}

\Fix{There is no such corrected pole-dominance transition.  The corrected leading pole satisfies
$\gamma_1<s_1<\alpha_1$ for all admissible $\beta>0$, so the asymptotic tail is controlled by
$s_1$ rather than by a competition between $s_1$ and $\alpha_1$.}

\Reason{The corrected finite partial-fraction expansion is
\begin{equation}
  F_\rho(t)=\sum_{j=1}^{L}B_{\rho j}s_j^{t-1},
\end{equation}
where the nonzero coefficients are residues:
\begin{equation}
  B_{\rho j}=\frac{qN_\rho(y_j)}{D'(y_j)},
  \quad
  N_\rho(y)=aT_{L-\rho}(y)+U_{\rho-1}(y)+U_{L-\rho-1}(y),
  \quad
  D(y)=aT_L(y)+U_{L-1}(y).
\end{equation}
Thus
\begin{equation}
  \Corrected{F_\rho(t)=B_{\rho 1}s_1^{t-1}+O(\sigma_2^{t-1}),
  \qquad \sigma_2=\max_{j\ge2}|s_j|<s_1.}
\end{equation}
The decay time is
\begin{equation}
  \Corrected{\tau_\beta=-\frac{1}{\log s_1}.}
\end{equation}}

\section{Audit item 6: perturbative dependence of the decay on beta}

For a referee asking specifically about the decay dependence of the second peak's right tail, the
most useful analytical dependence is obtained by perturbing the pole equation at weak shortcut
strength.  Let
\begin{equation}
  \lambda=\beta(1-q),
  \qquad a=\frac{q}{\lambda},
  \qquad \varepsilon=\frac{1}{a}=\frac{\lambda}{q}.
\end{equation}
Divide the pole equation by $a$:
\begin{equation}
  T_L(y)+\varepsilon U_{L-1}(y)=0.
\end{equation}
At $\varepsilon=0$ the largest root is
\begin{equation}
  y_0=\cos\frac{\pi}{2L}=\cos\frac{\pi}{N},
  \qquad
  s_0=1-q+qy_0=\alpha_1.
\end{equation}
Let $y=y_0+\varepsilon y_1+O(\varepsilon^2)$.  Substituting and retaining first order gives
\begin{equation}
  T_L'(y_0)y_1+U_{L-1}(y_0)=0,
  \qquad
  y_1=-\frac{U_{L-1}(y_0)}{T_L'(y_0)}.
\end{equation}
With $y_0=\cos\theta_0$, $\theta_0=\pi/(2L)$,
\begin{equation}
  U_{L-1}(y_0)=\frac{\sin(L\theta_0)}{\sin\theta_0}=\frac{1}{\sin\theta_0},
  \qquad
  T_L'(y_0)=L\frac{\sin(L\theta_0)}{\sin\theta_0}=\frac{L}{\sin\theta_0}.
\end{equation}
Hence $y_1=-1/L$ and
\begin{equation}
  s_1=1-q+qy=\alpha_1-q\frac{\lambda}{q}\frac{1}{L}+O(\lambda^2)
  =\alpha_1-\frac{\lambda}{L}+O(\lambda^2).
\end{equation}
Since $L=N/2$ and $\lambda=\beta(1-q)$,
\begin{equation}
  \Corrected{s_1=\alpha_1-\frac{2\beta(1-q)}{N}+O([\beta(1-q)]^2).}
\end{equation}
Therefore
\begin{equation}
  \Corrected{1-s_1=1-\alpha_1+\frac{2\beta(1-q)}{N}+O([\beta(1-q)]^2).}
\end{equation}
For large $N$,
\begin{equation}
  1-\alpha_1=q\left(1-\cos\frac{\pi}{N}\right)
  \sim \frac{q\pi^2}{2N^2},
\end{equation}
so
\begin{equation}
  \Corrected{1-s_1\sim \frac{q\pi^2}{2N^2}+\frac{2\beta(1-q)}{N}.}
\end{equation}
This is the requested decay-dependence item.

\section{Audit item 7: bimodality is not determined by tail dominance alone}

\Issue{The original text treats the long-time dominant eigenvalue comparison as an analytical
condition for the presence of a second peak.  That inference is too strong even aside from the sign
correction.}

\Fix{Use the exact finite-time expression to test local maxima.  If
\begin{equation}
  F_\rho(t)=\sum_{j=1}^{L}B_{\rho j}s_j^{t-1},
\end{equation}
then a local maximum at time $t_*$ requires
\begin{equation}
  \Delta F_\rho(t_*-1)>0,
  \qquad
  \Delta F_\rho(t_*)<0,
\end{equation}
where
\begin{equation}
  \Corrected{
  \Delta F_\rho(t)=F_\rho(t+1)-F_\rho(t)
  =\sum_{j=1}^{L}B_{\rho j}(s_j-1)s_j^{t-1}.}
\end{equation}
The existence of two sign changes in this discrete derivative is the direct analytical criterion
for two peaks.}

\Reason{The long-time tail answers how the right flank of a late peak decays if such a peak is
present.  It does not, by itself, prove that the finite-time curve has two local maxima.  The
correct referee response should separate these two questions: finite-time peak formation versus
right-tail decay.}

\section{Audit item 8: absorbing-ring convolution sign around the double sum}

\Issue{In the absorbing-ring section, the displayed double sum uses one bracket ordering, but the
next line effectively uses the opposite ordering.  The final compact formula can be made correct,
but the intermediate step must be aligned.}

\Fix{Let
\begin{equation}
  \lambda_j=1-q+q\cos\frac{\pi j}{N},
  \qquad
  \gamma_r=1-q+q\cos\frac{2\pi r}{N}.
\end{equation}
The time convolution
\begin{equation}
  \sum_{s=1}^{t} \lambda_j^{s-1}\gamma_r^{t-s}
\end{equation}
produces
\begin{equation}
  \frac{\gamma_r^t-\lambda_j^t}{\gamma_r-\lambda_j}
\end{equation}
when the denominator is written as $\gamma_r-\lambda_j$.  If one writes the denominator as
$\cos(\pi j/N)-\cos(2\pi r/N)$, the same expression must be arranged consistently.}

\Reason{This is a sign-bookkeeping issue rather than a conceptual failure of the absorbing-ring
calculation.  In the corrected derivation I spell out the convolution and keep the denominator and
numerator signs matched.}

\section{Audit item 9: initial-condition and boundary checks for the killed propagator}

\Issue{The original text gives oversimplified orthogonality statements after the compact killed
propagator.  They are not valid as written for all $n,n_0$.}

\Fix{The correct identities are, for $1\le d,d_0\le N-1$,
\begin{equation}
  \sum_{r=1}^{N-1}\cos\frac{2\pi rd}{N}\cos\frac{2\pi rd_0}{N}
  =\frac{N}{2}\bigl(\delta_{d,d_0}+\delta_{d,N-d_0}\bigr)-1,
\end{equation}
and
\begin{equation}
  \sum_{j=1}^{N-1}[1-(-1)^j]\sin\frac{\pi jd}{N}\sin\frac{\pi jd_0}{N}
  =\frac{N}{2}\bigl(\delta_{d,d_0}+\delta_{d,N-d_0}\bigr).
\end{equation}
Substituting these into the compact formula proves $W_{n_0}(n,0)=\delta_{n,n_0}$ for transient
sites.}

\Reason{The killed propagator also satisfies $W_{n_0}(m,t)=0$ and $W_m(n,t)=0$ once the convention
for an initially absorbing particle is fixed.  The original line says $Q$ where it should say $W$.}

\section{Audit item 10: first-passage generating function notation in the no-shortcut ring}

\Issue{The compact odd-mode notation is introduced for $k=1,
\ldots,N/2$ with
\begin{equation}
  \alpha_k=1-q+q\cos\frac{(2k-1)\pi}{N}.
\end{equation}
Immediately afterwards, the generating function is displayed with a sum $k=1,\ldots,N-1$ while
using the compact $f_k,\alpha_k$ notation.}

\Fix{Use either the full $j$-sum
\begin{equation}
  \wt F^{(0)}_{n_0}(m,z)=
  \frac{qz}{N}\sum_{j=1}^{N-1}
  \frac{[1-(-1)^j]\sin\bigl(|m-n_0|\pi j/N\bigr)\sin(\pi j/N)}
  {1-z[1-q+q\cos(\pi j/N)]},
\end{equation}
or the odd-mode $\ell$-sum
\begin{equation}
  \wt F^{(0)}_{n_0}(m,z)=z\sum_{\ell=1}^{N/2}
  \frac{f_\ell(n_0,m)}{1-z\alpha_\ell}.
\end{equation}}

\Reason{This avoids an ambiguity that becomes serious later when the symbols $\alpha_\ell$ and
$\gamma_r$ are compared in pole arguments.}

\section{Audit item 11: problems in the original convolution formula analogous to Eq. (41)}

\Issue{The original final finite-time convolution formula contains additional issues beyond the
shortcut sign.}

\Fix{At least the following should be corrected or removed in the revised manuscript:
\begin{enumerate}
  \item The no-shortcut first term should involve first passage to $v$, not to $u$:
  \begin{equation}
    \Original{f_\ell(n_0,u)\alpha_\ell^{t-1}}
    \quad\longrightarrow\quad
    \Corrected{f_\ell(n_0,v)\alpha_\ell^{t-1}}.
  \end{equation}
  \item In the odd-sector terms involving $g^{(2)}_r$, the associated eigenvalue is $\alpha_r$, not
  $\gamma_r$.  Any term where an odd-sector index is paired with $\gamma_r$ should be rederived.
  \item The $t\alpha_\ell^t$ term is a repeated-pole artifact of the original uncorrected
  convolution.  With the corrected denominator $aT_L+U_{L-1}$, there is no generic coincidence
  between the shortcut poles and the no-shortcut $\alpha_\ell$ poles; hence no generic
  $t\alpha_1^t$ long-time law.
\end{enumerate}}

\Reason{Because the entire Eq.-(41)-type expression is downstream of the wrong $H$ denominator, it
is safer to replace it with the corrected closed-form generating function and its partial-fraction
expansion rather than repairing the old convolution term by term.}

\section{Audit item 12: residue coefficient replacing the original $c_k$ expression}

\Issue{The old coefficient $c_k$ is tied to the denominator $a-\phi$, and the displayed formula has
apparent bracket and derivative typos.  It also depends on the old pole geometry.}

\Fix{The corrected residue coefficient follows directly from
\begin{equation}
  \wt F_\rho(z)=\frac{N_\rho(y)}{D(y)},
  \qquad
  y=1+\frac1q\left(\frac1z-1\right),
\end{equation}
where
\begin{equation}
  N_\rho(y)=aT_{L-\rho}(y)+U_{\rho-1}(y)+U_{L-\rho-1}(y),
  \qquad
  D(y)=aT_L(y)+U_{L-1}(y).
\end{equation}
Since $z=1/[1+q(y-1)]$, a simple pole at $y_j$ corresponds to
\begin{equation}
  s_j=1-q+qy_j.
\end{equation}
The coefficient of $s_j^{t-1}$ is
\begin{equation}
  \Corrected{B_{\rho j}=\frac{qN_\rho(y_j)}{D'(y_j)}.}
\end{equation}}

\Reason{This expression is shorter, less error-prone, and tied directly to the corrected pole
equation.}

\section{Audit item 13: figure caption and textual typos}

The following non-exhaustive typo list should be fixed in any revision:
\begin{enumerate}
  \item Document option \Original{idelinks} should be \Corrected{hidelinks}.
  \item The original source uses unavailable local style files; the revised files are standalone.
  \item A label such as \Original{\texttt{\textbackslash label\{eq:\}}} should be replaced by a
  meaningful label or removed.
  \item The notation \Original{$G^{(1)}_k(n_0,n,,m)$} should be
  \Corrected{$G^{(1)}_k(n_0,n,m)$}.
  \item The text ``when either $n=m$ or $n_0=m$, one has that $Q_{n_0}(n,t)=0$'' should use $W$, not
  $Q$.
  \item The Fig. 1 caption contains a logically impossible phrase equivalent to requiring $q>1$ for
  a nonzero shortcut probability.  The stochastic setting is $0<q<1$, with nonzero shortcut
  probability requiring $\beta(1-q)>0$.
  \item Several source brackets around $s_k-1$ in the old $c_k$ formula appear mismatched.
\end{enumerate}

\section{Alternative interpretations that should be ruled out explicitly}

The sign correction is forced only after the model is fixed.  The manuscript should rule out the
following alternatives:
\begin{enumerate}
  \item \textbf{Shortcut to target as killing.}  This is the corrected model used here.  The walker
  at $u$ jumps to absorbing $v$ with probability $\lambda=\beta(1-q)$ and exits the transient
  system.  Denominator: $1+z\lambda\wt W_u(u,z)$.
  \item \textbf{Probability-conserving shortcut on a nonabsorbing ring.}  A directed edge $u\to v$
  may be a rank-one probability-conserving perturbation before absorption is imposed.  If $v$ is
  later made absorbing, the transient reduction still creates killing at $u$ when the shortcut lands
  at $v$.
  \item \textbf{Reverse shortcut or removed killing.}  A formula with denominator
  $1-z\lambda\wt W_u(u,z)$ would correspond to a different perturbation, for example adding
  persistence in the transient state space rather than removing it.  It should not be described as
  a shortcut into an absorbing target.
\end{enumerate}

\section{Suggested referee-facing response}

A concise response to the referee's decay question could be:

\begin{quote}
We have rederived the symmetric shortcut-to-target problem using the transient-state transition
matrix.  The shortcut from $u$ to the absorbing target $v$ is an additional killing channel at $u$,
so the corrected resolvent contains the denominator $1+z\beta(1-q)\wt W_u(u,z)$ rather than
$1-z\beta(1-q)\wt W_u(u,z)$.  Consequently, the pole equation in the symmetric case $N=2L$ is
$aT_L(y)+U_{L-1}(y)=0$, with $a=q/[\beta(1-q)]$.  The right tail of the first-passage distribution,
including the right tail of a second peak if present, obeys
\begin{equation}
  F_\rho(t)\sim B_{\rho1}s_1^{t-1},
\end{equation}
where $s_1=1-q+q\cos\theta_1$ and $\theta_1$ solves
\begin{equation}
  \tan(L\theta_1)=-\frac{q}{\beta(1-q)}\sin\theta_1,
  \qquad
  \frac{\pi}{2L}<\theta_1<\frac{\pi}{L}.
\end{equation}
For weak shortcuts,
\begin{equation}
  1-s_1=1-\alpha_1+\frac{2\beta(1-q)}{N}+O([\beta(1-q)]^2)
  \sim \frac{q\pi^2}{2N^2}+\frac{2\beta(1-q)}{N}.
\end{equation}
Thus the earlier condition $\beta\le 2q/[(1-q)N]$ and the generic $t\alpha_1^t$ tail are not
properties of the corrected stochastic shortcut-to-target model.  Bimodality itself is a finite-time
shape property and should be checked from the exact finite difference of the corrected finite sum.
\end{quote}

\section{Checklist for the revised manuscript}

Before submitting a revised analytical supplement, check the following:
\begin{enumerate}
  \item Define the transition probabilities at $u$ explicitly.
  \item State whether the shortcut endpoint is absorbing at the instant of arrival.
  \item Replace the old shortcut resolvent by the killed rank-one resolvent.
  \item Replace the old $H$ denominator $a-\phi$ by $a+\phi$.
  \item Redraw or remove Fig. 1; if redrawn with the same $\phi(1/s)$ curve, use the line $-a$.
  \item Remove the threshold $\beta\le 2q/[(1-q)N]$ unless a separate finite-time bimodality proof
  establishes a new threshold.
  \item Replace the long-time statement by $F_\rho(t)\sim B_{\rho1}s_1^{t-1}$.
  \item Add the weak-shortcut decay-gap expansion.
  \item Separate the right-tail decay question from the finite-time two-peak-existence question.
  \item Use the exact finite difference $\Delta F(t)$ for any rigorous second-peak criterion.
\end{enumerate}

\end{document}
